\chapter{基于Vision-xLSTM局部--全局协同建模的三维医学图像分割方法}

% TODO(P-002): Work1 is a complete research loop. Use only paper-reported or approved source evidence.

\section{引言}

% TODO: Local 3D CNN strength -> long-range limitation -> 3D global-modeling cost -> ViL-UNet.

\section{ViL-UNet总体架构}

% TODO: Add a verified architecture figure; describe full data flow before modules.

\section{三维局部特征编码}

% TODO: Explain the actual role of 3D CNN in ViL-UNet, not generic CNN theory.

\section{Vision-xLSTM全局上下文建模}

\subsection{三维特征序列化与多方向建模}

% TODO: Recheck against the Work1 paper before formal writing.

\subsection{ViL模块及长距离信息交互}

% TODO: Recheck against the Work1 paper before formal writing.

\section{解码与特征融合}

% TODO: Describe only verifiable structures in the Work1 paper.

\section{实验与结果分析}

\subsection{数据集与实验设置}

% TODO: Verify Synapse and ACDC protocol details from reliable Work1 sources.

\subsection{评价指标与对比方法}

% TODO: Refer to Chapter 2 definitions; record the actual metrics and baselines used.

\subsection{Synapse实验结果与分析}

% TODO: PAPER_REPORTED results only; obtain organ-level data before adding it.

\subsection{ACDC实验结果与分析}

% TODO: PAPER_REPORTED results only.

\subsection{ViL模块数量消融实验}

% TODO: Analyze the reported 3/6/12/18/24-block ablation with cautious causal language.

\subsection{模型复杂度与效率分析}

% TODO: PAPER_REPORTED Params and FLOPs only; do not add unverified VRAM or inference time.

\subsection{可视化分析}

% TODO: Retain only if an authentic Work1 figure or later lawful experiment output is available.

\section{本章小结}

% TODO: Problem -> ViL-UNet -> core design -> key evidence -> limited conclusion; bridge to scale-level structure.
